%!TEX program = xelatex
\documentclass[12pt, b5paper]{article}
\usepackage{ctex}

\usepackage{geometry}
\geometry{b5paper,left=2cm,right=1cm,top=2.5cm,bottom=2.5cm}
\usepackage{chemformula} %化学公式宏包
\usepackage[version=4]{mhchem} %化学公式宏包
\usepackage{siunitx} %数字，单位宏包
\sisetup{
	separate-uncertainty = true,
	inter-unit-product = \ensuremath{{}\cdot{}}
} %单位间加小圆点
\usepackage{tikz} %Tikz绘图包
\usetikzlibrary{cd}
\usetikzlibrary{chains}
\usetikzlibrary{arrows}
\usetikzlibrary{arrows.meta} %画箭头
\usetikzlibrary{commutative-diagrams} %交换图
\usetikzlibrary{positioning,automata}
\usepackage[all]{xy}
\usepackage{booktabs} %三线表
\usepackage{multirow} %合并单元格
\usepackage{makecell} %表格内强制换行
\usepackage{array}
\usepackage{booktabs} %控制表格行高
\usepackage{colortbl}  %彩色表格需要加载的宏包
\usepackage{amsmath}
\usepackage{unicode-math}
\usepackage{fontspec} %字体
\newcommand*{\cred}{\textcolor{red}}

\setmainfont{Times New Roman}
\setmathfont{XITS Math}

\setlength{\parindent}{0pt} %无缩进
\usepackage{setspace}
\usepackage{fancyhdr} %设置页码
\usepackage{lastpage} %获取总页数
\pagestyle{fancy} %fancyhdr宏包新增的页面风格
\renewcommand\headrulewidth{0pt} %取消页眉中的装饰分割线
\fancyhf{}
\cfoot{\footnotesize 第 \thepage\ 页(共 \pageref{LastPage}页)}%当前页 of 总页数

\begin{document}
\begin{spacing}{1.625}

\centerline{{\Large 河北农业大学2022——2023学年第1学期}}

\centerline{\Large 参考答案及评分标准（A）卷}

\centerline{开课学院：\underline{\quad 理工学院 \quad} \quad 出\ \ 题\ \ 人：\underline{ \qquad \qquad 张斌 \qquad \qquad}}

\centerline{课\ \ 程\ \ 号：\underline{HBX250037-01} \quad 课程名称：\underline{\quad \quad \ \ 化学反应工程 \quad \ \ \ }}

\bigskip


1. 简答题(15分)


要求能够将反应工程某一章节的主要内容知识点总结清楚。
\bigskip
\bigskip

2. 简答题(15分)

要求针对某一具体问题进行分析，将解决问题中遇到的困难，以及是如何解决困难的进行说明。

\bigskip
\bigskip

3. 简答题(16分)

可以从各种类型反应器的特点、操作方式、反应器计算等方面进行回答。

\bigskip
\bigskip

4. 读图题(24分)

(1) 可逆放热反应 $\hfill \qquad \cdots \cdots \cdots \cdots (4\text{分})$

(2) M点或N点，视情况而定。$\hfill \qquad \cdots \cdots \cdots \cdots (4\text{分})$

(3) MA等温操作、GD绝热操作、GK换热操作。$\hfill \qquad \cdots \cdots \cdots \cdots (6\text{分})$

(4) 热点。$\hfill \qquad \cdots \cdots \cdots \cdots (4\text{分})$

反应热效应、床层传热强度、进料温度、操作变量等。$\hfill \qquad \cdots \cdots \cdots \cdots (6\text{分})$

\bigskip
\bigskip

5. 计算题(30分)

设小釜体积为$V_{r1}$，大釜体积为$V_{r2}$

（1）大釜在前，小釜在后：

$\begin{cases}
    \frac{V_{r2}}{Q_0} = \frac{c_{A0}X_{A1}}{kc_{A0}(1-X_{A1})} \\
    \frac{V_{r1}}{Q_0} = \frac{c_{A0}(X_{A2} - X_{A1})}{kc_{A0}(1 - X_{A2})}
\end{cases}$

出口转化率$X_{A2} = \frac{kV_{r2} + \frac{kQ_0V_{r1}}{Q_0 + kV_{r1}}}{Q_0 + kV_{r2}}\hfill \qquad \cdots \cdots \cdots \cdots (10\text{分})$ 

（2）小釜在前，大釜在后：

$\begin{cases}
    \frac{V_{r1}}{Q_0} = \frac{c_{A0}X'_{A1}}{kc_{A0}(1-X'_{A1})} \\
    \frac{V_{r2}}{Q_0} = \frac{c_{A0}(X'_{A2} - X'_{A1})}{kc_{A0}(1 - X'_{A2})}
\end{cases}$

出口转化率$X'_{A2} = \frac{kV_{r1} + \frac{kQ_0V_{r2}}{Q_0 + kV_{r2}}}{Q_0 + kV_{r1}}\hfill \qquad \cdots \cdots \cdots \cdots (10\text{分})$

$X_{A2} = X'_{A2} = \frac{kQ_0V_{r2} + kQ_0V_{r1} + k^2V_{r1}V_{r2}}{(Q_0 + kV_{r2})(Q_0 + kV_{r1})}$

所以出口最终转化率相同。$\hfill \qquad \cdots \cdots \cdots \cdots (10\text{分})$

代入特定数据计算扣10分。


\end{spacing}
\end{document}
